\documentclass[11pt]{article}
\usepackage{amsmath,amssymb,amsthm}
\setlength{\emergencystretch}{2em}
\newtheorem{theorem}{Theorem}
\newtheorem{lemma}[theorem]{Lemma}
\newtheorem{proposition}[theorem]{Proposition}
\newtheorem{corollary}[theorem]{Corollary}
\newcommand{\im}{\operatorname{im}}
\newcommand{\Sp}{\operatorname{Sp}}
\newcommand{\End}{\operatorname{End}}
\newcommand{\Fix}{\operatorname{Fix}}
\begin{document}
\title{Characteristic vectors and endpoint geometry}
\author{Human Compression}
\date{}
\maketitle

We study involutions in a six-dimensional symplectic space in
characteristic two. Our main result excludes an irreducible string
marking of rank at least five when both endpoints have maximal rank
after subtracting the identity. The proof uses at most three translates
of one characteristic vector. It requires neither a second-generator
centralizer nor an enumeration of completions.

The characteristic-vector construction is borrowed from the existing
project note \emph{Canonical fixed vectors in characteristic two},
Theorem~1, at
\texttt{\detokenize{research/groups/common_commutant/fixed_vectors.tex}}
in the October corpus. The Wall forms, adjoints, rank bounds and
Cayley--Hamilton identity used below are standard linear algebra.
The additional argument here combines these facts with the string
relations. No assertion of historical priority is made.
Applications to a named finite group still require the original group,
representation, endpoint and ambient-escape premises. No application
to an unacquired endpoint pair, census result, conjecture solution,
formal verification or numerical compression is asserted.

Except in the explicitly labelled source example, vectors are
\emph{columns}. Let \(k\) be a perfect field of characteristic two,
let \(V\) have dimension six, and let \(\beta\) be a nondegenerate
alternating form on \(V\). A string marking
\((t_0,\ldots,t_{r-1})\) consists of involutions in \(\Sp(V,\beta)\)
satisfying \([t_i,t_j]=1\) for \(|i-j|\geq2\).
No intersection property is assumed.

\section{Characteristic vectors and Wall forms}

For an involution \(s\in\Sp(V,\beta)\), put \(N_s=s-I\) and
\(P_s=\im N_s\). The operator \(N_s\) is self-adjoint and
\(N_s^2=0\). The borrowed characteristic-vector theorem gives a
unique vector \(c_s\) such that
\[
             \beta(x,c_s)^2=\beta(x,N_sx)
                         \qquad(x\in V).                 \tag{1}
\]
It belongs to \(P_s\). Every symplectic isometry commuting with \(s\)
fixes \(c_s\) pointwise. The vector can be zero.

The Wall form on \(P_s\) is
\[
        h_s(N_sx,N_sy)=\beta(x,N_sy).                     \tag{2}
\]
It is well defined: changing \(x\) by an element of \(\ker N_s\)
does not change its pairing with \(\im N_s\), since
\((\im N_s)^\perp=\ker N_s\). Self-adjointness makes it symmetric.
If \(h_s(N_sx,P_s)=0\), then
\(\beta(x,P_s)=0\), so \(x\in\ker N_s\) and \(N_sx=0\).
Thus \(h_s\) is nondegenerate.
Equations (1) and (2) give
\[
             h_s(u,c_s)^2=h_s(u,u)\qquad(u\in P_s).       \tag{3}
\]
In particular, \(c_s=0\) precisely when \(h_s\) is alternating.
Every symplectic isometry commuting with \(s\) induces an isometry
of \(h_s\): if \(u=N_sx\), it sends this representative to
\(N_sgx=gN_sx\), and preservation of (2) follows from that of
\(\beta\).

\begin{lemma}\label{lem:oddwall}
Suppose that \(\operatorname{rank}N_s=3\). Then
\[
 P_s=\ker N_s=P_s^\perp,\qquad
 c_s\ne0,\qquad h_s(c_s,c_s)=1.
\]
Consequently
\[
       P_s=kc_s\mathbin{\perp_{h_s}} W_s,\qquad
       W_s=c_s^{\perp_{h_s}},
                                                               \tag{4}
\]
where \(W_s\) is a two-dimensional symplectic space for \(h_s\).
\end{lemma}
\begin{proof}
The image and kernel of \(N_s\) have dimension three, and its square
is zero, so they coincide. Self-adjointness gives
\(P_s^\perp=\ker N_s\). A nondegenerate alternating form cannot
have odd dimension, so the nondegenerate symmetric form \(h_s\)
on \(P_s\) is not alternating. Equation (3) implies \(c_s\ne0\).
Putting \(u=c_s\) in (3) shows that \(h_s(c_s,c_s)\) is either
zero or one.

If it were zero, \(c_s\) would lie in the two-dimensional space
\(c_s^{\perp_{h_s}}\). Equation (3) says that \(h_s\) restricts
to an alternating form on that space. Its radical would be exactly
\(kc_s\), because the orthogonal complement in \(P_s\) of that
plane is \(kc_s\). This would be an alternating form of rank one,
which is impossible. Hence \(h_s(c_s,c_s)=1\).
Its orthogonal complement is therefore nondegenerate, and (3)
makes that restriction alternating. This proves (4).
\end{proof}

\begin{lemma}\label{lem:commutingmaximal}
Let \(A,E\) be commuting symplectic involutions, each of rank three
after subtracting the identity. Put
\[
 N=A-I,\quad M=E-I,\quad P=\im N,\quad U=\im M,\quad
 H=P\cap U.
\]
Then \(\dim H\) is two or three. If it is two, then
\[
 L=\im(NM)
\]
is a line, and
\[
 H=kc_A+L=kc_E+L,\qquad c_A\notin L,\quad c_E\notin L.  \tag{5}
\]
The line \(L\) is isotropic for both \(h_A\) and \(h_E\).
\end{lemma}
\begin{proof}
By Lemma~\ref{lem:oddwall}, \(P=\ker N\) and \(U=\ker M\).
Commutation implies that \(N\) preserves \(U\), with kernel
\(P\cap U\). Writing \(d=\dim H\), we obtain
\[
       \operatorname{rank}(NM)=3-d.
\]
Its image lies in \(H\), so \(3-d\leq d\), whence \(d\geq2\).
If \(d=2\), this image is a line \(L\).

The restriction of \(E\) to \(P\) is an involutory isometry of
\(h_A\). Thus \(M|_P\) is self-adjoint for \(h_A\), has square
zero, and has image \(L\). For \(u,v\in P\),
\[
 h_A(Mu,Mv)=h_A(u,M^2v)=0.
\]
Hence \(L\) is isotropic for \(h_A\). Since
\(h_A(c_A,c_A)=1\), it does not contain \(c_A\).
The vector \(c_A\) is fixed by \(E\), because \(E\) commutes with
\(A\); thus \(c_A\in H\). Dimension two now gives
\(H=kc_A+L\). Interchanging \(A,E\) proves explicitly that
\(L\) is isotropic for \(h_E\), that \(c_E\notin L\), and that
\(H=kc_E+L\).
\end{proof}

\section{Two maximal endpoints}

\begin{theorem}\label{thm:maximal}
Let \((t_0,\ldots,t_{r-1})\) be a string marking in
\(\Sp(V,\beta)\), with \(r\geq5\). Suppose that
\[
       \operatorname{rank}(t_0-I)
        =\operatorname{rank}(t_{r-1}-I)=3.              \tag{6}
\]
Write \(A=t_0\), \(E=t_{r-1}\), \(P=\im(A-I)\) and
\(U=\im(E-I)\).

If \(P=U\), the generated subgroup stabilizes that Lagrangian.
If \(P\ne U\), the generated subgroup fixes either \(c_A\)
or \(c_E\), both of which are nonzero.
In particular, the generated subgroup is reducible on \(V\).
\end{theorem}
\begin{proof}
Every mark commutes with at least one endpoint. If \(P=U\), it
therefore preserves that space.

Suppose \(P\ne U\). The endpoints commute by the string relation.
Lemma~\ref{lem:commutingmaximal} applies with \(\dim(P\cap U)=2\)
and line \(L=\im((A-I)(E-I))\). For the Wall form \(h_E\), put
\[
              c=c_E,\qquad W=c^{\perp_{h_E}}\subset U.
\]
Lemma~\ref{lem:oddwall} gives \(U=kc\perp W\), with \(W\)
symplectic of dimension two.
All elements centralizing \(E\) preserve this decomposition.

On \(U\), the difference \(A-I\) has rank one and image \(L\).
Since it is self-adjoint for \(h_E\) and kills \(c\), we have
\(L\subset W\). Its restriction to \(W\) is nontrivial: otherwise
it would vanish on both summands of \(U\).
Thus the restriction of \(A\) to \(W\) is a nontrivial
involutory isometry whose difference from the identity has image
and kernel \(L\).
By (5), for any nonzero \(\ell\in L\) we can write
\[
                       c_A=a c+b\ell,\qquad a\ne0.     \tag{7}
\]
The inequality follows because \(c_A\notin L\).

Write \(B=t_1\). Every mark except \(B\) commutes with \(A\)
and therefore fixes \(c_A\).
If \(Bc_A=c_A\), the entire marking fixes \(c_A\), proving
the first alternative.

Suppose now \(Bc_A\ne c_A\). As \(B\) commutes with \(E\),
it fixes \(c\), preserves \(W\), and acts there as a
nontrivial involutory symplectic isometry.
Its difference from the identity on \(W\) has a one-dimensional
image and kernel, say \(L_B\).
If \(L_B=L\), it fixes \(\ell\) and \(c\), and hence fixes \(c_A\)
by (7), contrary to assumption. Therefore \(L_B\ne L\).

Let \(R=(AB)|_W\). It has no nonzero fixed vector.
Indeed, \(ABw=w\) implies \(Aw=Bw\), so
\((A-I)w=(B-I)w\) belongs to \(L\cap L_B=0\).
This forces \(w\in L\cap L_B=0\).
Since \(\det R=1\), its characteristic polynomial is
\[
                        X^2+\tau X+1.
\]
The absence of a fixed vector implies \(\tau\ne0\):
\(\det(R-I)=\tau\).
The Cayley--Hamilton identity gives
\[
                    R^2+\tau R+I=0\quad\hbox{on }W.
\]
The operator \(g=AB\) fixes \(c\), so the same polynomial is
multiplication by \(\tau\) on \(kc\). Applying it to (7) yields
\[
               (g^2+\tau g+I)c_A=a\tau c.             \tag{8}
\]

The penultimate mark \(D=t_{r-2}\) commutes with \(A\) and \(B\),
because \(r\geq5\). It fixes \(c_A\) and commutes with \(g\).
It therefore fixes the left-hand side of (8).
Since \(a\tau\ne0\), it fixes \(c=c_E\).
Every other mark either equals \(E\) or commutes with \(E\);
these marks fix \(c_E\) as well.
This proves the second alternative.
\end{proof}

\begin{corollary}\label{cor:originalgroup}
Let \(\rho:G\to\Sp(V,\beta)\) be a representation, and fix
endpoint elements \(a,e\in G\) whose images satisfy (6).
If those images have different images after subtracting the identity,
assume that \(\rho(G)\) moves each of \(c_A,c_E\).
If those images coincide in a space \(P\), assume instead that
\(\rho(G)\) does not stabilize \(P\).
Then no string marking of \(G\) of rank at least five with endpoints
\(a,e\) generates \(G\).
\end{corollary}
\begin{proof}
Theorem~\ref{thm:maximal} puts the image of the marking in a
vector or subspace stabilizer missing an element of \(\rho(G)\).
If the marking generated \(G\), its image would generate \(\rho(G)\).
Faithfulness is unnecessary.
\end{proof}

For different endpoint images, the two required vector escapes may
be different ambient elements. An irreducibility premise would imply
all the needed escapes, but explicit acquired escapes are preferable
when certifying a particular finite-group application.

\begin{corollary}\label{cor:oneambient}
In Corollary~\ref{cor:originalgroup}, suppose that
\[
                     (A-I)(E-I)\ne0.
\]
If \(\rho(G)\) contains an element \(d\) for which \(d-I\) is
invertible, then no such marking generates \(G\).
\end{corollary}
\begin{proof}
Lemma~\ref{lem:commutingmaximal} shows that the product is zero
exactly when \(P=U\). Thus Theorem~\ref{thm:maximal} gives a
nonzero vector fixed by the marking. The element \(d\) fixes no
nonzero vector, so cannot belong to its generated image.
\end{proof}

For example, in the previously retained exponent codec the diagonal
ambient matrix
\[
 \operatorname{diag}(\text{code }2,\text{code }7,\text{code }3,
                     \text{code }6,\text{code }2,\text{code }7)
\]
has no entry equal to code \(1\), and therefore its difference
from the identity is invertible. If its original ambient-member
correspondence and the same identified model apply to a proposed
endpoint family, it supplies this escape without computing either
characteristic vector. The displayed shape alone does not establish
that original correspondence.

The theorem does not identify which retained endpoint families satisfy
(6). No coverage of the remaining original jobs follows until their
original pair, representation and rank premises have been bound.

\paragraph{The residual endpoint geometry is nonempty.}
Lemma~\ref{lem:commutingmaximal} does not silently eliminate its
independent-vector case. Here is an explicit pair realizing it.
On \(k^4\) use the alternating Gram matrix
\(J_4=I_4+\mathbf1\mathbf1^{\mathsf T}\), where
\(\mathbf1=(1,1,1,1)^{\mathsf T}\).
It is nondegenerate because \(J_4^2=I_4\).
Let \(A_4,E_4\) be the permutation matrices of
\((12)(34)\), \((13)(24)\), respectively.
They commute, their differences from the identity have rank two,
and their characteristic vectors are both \(\mathbf1\).
Their image intersection is \(k\mathbf1\), and
\((A_4-I)(E_4-I)=\mathbf1\mathbf1^{\mathsf T}\).

On a separate symplectic plane with basis \(e,f\), let
\(T(f)=e\), \(T(e)=0\), and choose distinct nonzero \(a,b\in k\).
Set
\[
 A=A_4\oplus(I+aT),\qquad E=E_4\oplus(I+bT).
\]
The endpoints commute and each difference has rank three.
Their image intersection is
\(k(\mathbf1,0)+k(0,e)\), and
\[
 c_A=(\mathbf1,\sqrt a\,e),\qquad
 c_E=(\mathbf1,\sqrt b\,e)
\]
are independent.
All these assertions follow directly from the displayed blocks:
for example \(J_4(A_4-I)=A_4-I\) has diagonal all ones,
and \((aT)(bT)=0\).
This is an abstract symplectic example, not an acquired pair in the
original \(G_2(8)\) action. The new step in
Theorem~\ref{thm:maximal} is the orbit argument (8), not a claim
that independent characteristic vectors cannot occur.

\section{An algebraic witness for a prefix}

For arbitrary involutions \(A,B\in\Sp(V,\beta)\), write
\(\mathcal A=k[A,B]\). For each self-adjoint \(P\in\mathcal A\),
let \(w(P)\) be the borrowed characteristic vector satisfying
\[
                    \beta(x,w(P))^2=\beta(x,Px).
\]
Define
\[
 \mathcal W(A,B)=\{w(P):P\in\mathcal A,\ P=P^*\}.
                                                               \tag{9}
\]

\begin{proposition}\label{prop:prefix}
The set \(\mathcal W(A,B)\) is a \(k\)-subspace fixed pointwise
by \(C_{\Sp(V,\beta)}(A,B)\).
Let a string marking of rank at least five begin with \(A,B\)
and end with \(E\). If
\[
                       0\ne c_E\in\mathcal W(A,B),     \tag{10}
\]
then the whole marking fixes \(c_E\).
\end{proposition}
\begin{proof}
The characteristic-vector definition and uniqueness give
\[
 w(\alpha^2P+\gamma^2Q)=\alpha w(P)+\gamma w(Q).
\]
Because \(k\) is perfect, this proves that (9) is a subspace.
Every symplectic isometry commuting with \(A,B\) commutes with
every member of \(\mathcal A\), and therefore fixes each \(w(P)\).
The penultimate mark centralizes \(A,B\), so it fixes \(c_E\)
under (10). Every other mark centralizes \(E\) or equals \(E\),
and hence also fixes \(c_E\).
\end{proof}

This use of the borrowed theorem requires no equations saying that
other marks fix \(\im P\): their commutation with \(E\) already
fixes the particular vector \(c_E\). No generators for the full
endpoint centralizer are needed.

Write \(\beta(x,y)=x^{\mathsf T}Jy\). A sufficient exact witness
for (10) consists of a word expression
\[
                         P=\sum_w\lambda_w w(A,B)
\]
and the two identities
\[
        P^{\mathsf T}J=JP,\qquad
        \operatorname{diag}(JP)=(Jc_E)^{[2]}.           \tag{11}
\]
Here brackets denote entrywise squaring. The borrowed formula
\(w(P)=J^{-1}\sqrt{\operatorname{diag}(JP)}\) proves the claim.
Membership in the algebra is certified by the displayed expression;
completeness of an algebra basis is not needed.
For a fixed spanning list of words, (11) is an inhomogeneous
linear system in their coefficients. A failed search in an
incomplete list proves nothing.

In the moving case of Theorem~\ref{thm:maximal}, set \(g=AB\) and
\[
       P_i=g^i(A-I)g^{-i}\quad(i=0,1,2),\qquad
       P=\frac{P_2+\tau^2P_1+P_0}{(a\tau)^2}.         \tag{12}
\]
Each \(P_i\) is self-adjoint, belongs to \(k[A,B]\), and has
\(w(P_i)=g^ic_A\). Equation (8) therefore gives \(w(P)=c_E\).
Thus (12) is a uniform explicit witness using three conjugates.
The scalars \(a=h_E(c_A,c_E)\) and
\(\tau=\operatorname{tr}((AB)|_{W_E})\) are nonzero in the moving branch
where that formula is used.

Two simpler sufficient conditions also follow from (9).
First, \(c_E\in\langle c_A,c_B\rangle\setminus\{0\}\) implies
(10), using a linear combination of \(A-I,B-I\).
Second, \(\mathcal W(A,B)\) contains the span of all vectors
\[
       g c_A,\quad g c_B\qquad(g\in\langle A,B\rangle),
\]
because conjugation transports characteristic vectors.
These are sufficient criteria, not asserted descriptions of all
of \(\mathcal W(A,B)\).

\section{What changes for mixed ranks}

\begin{proposition}\label{prop:mixed}
Let \(A,E\) be commuting symplectic involutions with
\[
 \operatorname{rank}(A-I)=2,\quad
 \operatorname{rank}(E-I)=3,\quad (A-I)(E-I)\ne0.
\]
Put \(P=\im(A-I)\), \(U=\im(E-I)\), \(c=c_E\), and
\(L=\im((A-I)(E-I))\). Then
\[
 \dim L=1,\qquad P\cap U=L,\qquad c\notin P,
\]
and
\[
                       Z=P+kc
\]
is a Lagrangian distinct from \(U\), with
\[
 U\cap Z=L+kc.
\]
The involution \(A\) fixes \(Z\) pointwise and has rank one on
\(U\) after subtracting the identity. The involution \(E\)
fixes \(U\) pointwise and has rank one on \(Z\) after subtracting
the identity.
\end{proposition}
\begin{proof}
On the three-dimensional \(U\), the involution \(A\) has
nonzero square-zero difference, hence rank one. Its image is \(L\).
It is an isometry of \(h_E\), so \(L\) is isotropic for \(h_E\).
Lemma~\ref{lem:oddwall} gives \(h_E(c,c)=1\); thus \(c\notin L\).

The plane \(P\) cannot be contained in \(U=\ker(E-I)\),
since \((E-I)P=\im((E-I)(A-I))=L\ne0\).
As \(L\subseteq P\cap U\), their intersection is exactly \(L\).
In particular \(c\notin P\).
Commutation with \(E\) makes \(A\) fix \(c\), so
\(c\in\ker(A-I)=P^\perp\).
The plane \(P\) is isotropic, and therefore its extension
\(P+kc\) is an isotropic three-space, hence Lagrangian.
The stated intersection and actions now follow from the definitions.
\end{proof}

This gives a canonical configuration of two Lagrangians with a
common plane. It does not say that the intermediate marks preserve
either of them. In this case the universal theorem from the separate
note \emph{Symplectic flags in string generating sets}, whose
product hypothesis is \((A-I)(E-I)=0\), cannot simply be extended
by deleting that hypothesis.

\paragraph{The original positive construction gives a boundary.}
The following example uses \emph{row vectors throughout this paragraph}.
In Chen's source edition \cite{Chen}, take an even prime power \(q>2\),
a generator \(\omega\in\mathbb F_{q^2}^{\,*}\), and
\[
  \zeta=\omega^{q-1},\qquad \lambda=\zeta+\zeta^{-1}\in\mathbb F_q.
\]
Here \(\zeta\) has order \(q+1\), and \(\lambda\ne0,1\).
These are the source's Section~3 hypotheses; no arbitrary replacement
of the parameter is intended.
The source's seven-dimensional model fixes \(y_4\).
On the quotient, use row basis
\[
                   (y_1,y_2,y_3,y_5,y_6,y_7)
\]
and the antidiagonal symplectic form.
The first and last displayed source generators become
\[
 A=t_1=I+\lambda(E_{23}+E_{45}),\qquad
 E=t_5=I+E_{16}+E_{24}+E_{35}.                        \tag{13}
\]
With the row convention, \(P\) denotes the row image of \(A-I\).
Direct symbolic multiplication gives
\[
  \operatorname{rank}(A-I)=2,\quad
  \operatorname{rank}(E-I)=3,\quad
  (A-I)(E-I)=(E-I)(A-I)=\lambda E_{25}\ne0.
\]
For a row \(x\),
\[
 \beta(x,x(A-I))=0,\qquad \beta(x,x(E-I))=x_1^2.
\]
Thus \(h_A\) is alternating and \(c_E=e_6\), while
\[
 P=\langle e_3,e_5\rangle,\quad
 U=\langle e_4,e_5,e_6\rangle,\quad
 Z=\langle e_3,e_5,e_6\rangle.
\]
The second source generator swaps \(e_2,e_3\) and \(e_4,e_5\).
It moves \(P\) and \(Z\).
The fourth source generator is the permutation matrix for
\((1\,6)(2\,4)(3\,5)\); it moves \(U\) and \(Z\).
These facts follow directly from the source's displayed matrices,
without any enumeration.

The assertion that the five original generators form a C-string
generating \(G_2(q)\) is imported from the source's Theorem~1.1.
It is not reproved here. In particular this source construction
precludes an indiscriminate nongeneration assertion for all mixed
rank-two/rank-three commuting endpoints with nonzero product.
It also directly shows that neither the old pair of subspaces
\(P,U\), nor the new mixed flag \(Z\), must be invariant under
the whole marking.

To compare with the column convention elsewhere in this note,
take the contragredient representation
\(g\mapsto \rho_{\rm row}(g^{-1})^{\mathsf T}\).
For the involution marks this replaces each displayed matrix by
its transpose. It preserves the relevant ranks, commutation
relations and zero/nonzero product distinctions.
The row and column subspaces must be transported accordingly;
their coordinate names are not interchangeable.

\section{Premises still needed for finite applications}

\begin{lemma}\label{lem:projective}
Let \(q\) be even, let \(V\) have dimension \(n\) over
\(\mathbb F_q\), and let \(s\in\operatorname{GL}(V)\) satisfy
\(s^2=I\). In its action on \emph{all} one-dimensional subspaces
of \(V\), the number of fixed points is
\[
             \frac{q^{\,n-\operatorname{rank}(s-I)}-1}{q-1}.
                                                               \tag{14}
\]
\end{lemma}
\begin{proof}
If \(s\) preserves a line \(kx\), then \(sx=\lambda x\) for
some \(\lambda\ne0\). The relation \(s^2=I\) gives
\(\lambda^2=1\), whence \(\lambda=1\) in characteristic two.
Thus its fixed projective points are exactly the lines in
\(\ker(s-I)\). A \(d\)-dimensional space has \(q^d-1\) nonzero
vectors, with exactly \(q-1\) on each line. Rank-nullity gives
(14).
\end{proof}

For \(n=6,q=8\), the full projective degree is \(37449\).
Formula (14) gives \(73\) fixed points for rank three and
\(585\) for rank two. Conversely, these counts determine those
ranks, because \((8^d-1)/7\) is strictly increasing in \(d\).
For two commuting symplectic involutions of rank three,
Lemma~\ref{lem:commutingmaximal} further gives:
their common fixed projective points number \(73\) when their
image Lagrangians coincide, and \(9\) when those Lagrangians
are distinct.
Alternatively, once both individual counts are \(73\), a single
projective point fixed by one endpoint and not the other proves
that the Lagrangians are distinct and that their nilpotent product
is nonzero.

These observations offer a possible smaller premise certificate,
not an acquired one. To use them with an original permutation
action on \(\Omega\), one must retain an explicit equivariant
identification
\[
                      \Omega\longleftrightarrow\mathbb P(V)
\]
for the same original group and representation.
The numerical degree \(37449\), an orbit of that size, or the
names of permutation files do not establish this identification.
No actual fixed-point counts or point correspondences for the
remaining original endpoint families are asserted here.

For a fixed original pair, the maximal-endpoint theorem requires:
an original representation into the stated symplectic group in a
single identified basis and field codec; exact original endpoint
images of rank three; their ordering in every covered original
marking; and the relevant acquired ambient vector or subspace escapes.
The string relations supply endpoint commutation.
It does not require acquiring a second generator, a full group
centralizer, a linear commutant or any completed marking.

For the more general algebraic criterion, one additionally needs the
actual original second-generator image and a word expression satisfying
(11), together with an ambient element moving \(c_E\).
The original model, membership, transport and field statements cannot
be replaced by a coincidental matrix shape.
Semilinear outer automorphisms require their existing original
transport argument; they are not automatically matrices in this
six-dimensional linear group.

Nothing here proves that all remaining \(G_2(8)\) endpoint families
have maximal-rank endpoints, or that the mixed cases satisfy (10).
No original job is discharged by this note alone.

\begin{thebibliography}{1}
\bibitem{Chen}
M. H. W. Chen,
\emph{Abstract regular polytopes of \(G_2(q)\) in even characteristic},
arXiv:2609.25004v1, 21 September 2026.
The retained extracted source is
\texttt{\detokenize{research/leemans/sources/g2_even2026.txt}},
SHA256
\texttt{\detokenize{a856dbe5a29e96b7545c4ab295a31969f9306b9079e38a74a3365590fe0bbf9c}}.
\end{thebibliography}
\end{document}
